%----------------------------------------------------------------------------------------
% NOTATIE EN FIGUREN VAN DE CURSUS LOGICA (2de graad)
% Staat in de preambule en niet in een hoofdstuk: ook de lijsten met oplossingen gebruiken dit.
% Vertrokken van notatie.tex van Integralen; de grafieken (rooster, kromme, gebied ...) zijn geschrapt,
% de waarheidstabellen en de logische symbolen zijn erbij gekomen.
%----------------------------------------------------------------------------------------
% Bewaarde grafiek, geplaatst door maak-pdf.py: bestand en pagina in figuren-cache/
\newcommand{\grafiekbestand}[2]{\includegraphics[page=#2,trim=3bp 3bp 3bp 3bp]{figuren-cache/#1}}

% Zoals stappen, maar met een bredere kolom voor de berekening: voor lange formules
\newenvironment{stappenbreed}{%
  \par\renewcommand{\arraystretch}{1.25}%
  \noindent\begin{tabular}{@{}>{$\displaystyle}p{9.8cm}<{$}@{\hspace{4mm}}>{\RaggedRight\small\itshape\arraybackslash}p{4.9cm}@{}}%
}{\end{tabular}\par}

% Zoals stappen, met eigen kolombreedtes: \begin{stappenvrij}{6.2cm}{3.8cm}. Voor naast een figuur.
\newenvironment{stappenvrij}[2]{%
  \par\renewcommand{\arraystretch}{1.25}%
  \noindent\begin{tabular}{@{}>{$\displaystyle}p{#1}<{$}@{\hspace{3mm}}>{\RaggedRight\small\itshape\arraybackslash}p{#2}@{}}%
}{\end{tabular}\par}

% Tekst of berekening links, figuur rechts: \naast[breedte rechts]{links}{rechts}
% De breedte rechts komt uit de figuur zelf; het optionele argument wordt niet gebruikt.
\newsavebox{\naastbox}
\newcommand{\naast}[3][]{%
  \par\noindent\sbox{\naastbox}{#3}%
  \begin{minipage}[t]{\dimexpr\linewidth-\wd\naastbox-3mm\relax}\vspace{0pt}#2\end{minipage}\hfill
  \begin{minipage}[t]{\wd\naastbox}\vspace{0pt}\usebox{\naastbox}\end{minipage}\par\addvspace{1mm}}
% Een rij figuren, gecentreerd
\newcommand{\figuren}[1]{\par\addvspace{1mm}{\centering #1\par}\addvspace{1.5mm}}
% Opgave met een grafiek rechts: \begin{metgrafiek}[breedte]{tikz} opgave en deelvragen \end{metgrafiek}
% De solution komt na \end{metgrafiek}, zoals bij metprent.
\newsavebox{\grafiekbox}
\newenvironment{metgrafiek}[2][]{%
  \sbox{\grafiekbox}{#2}%
  \begin{minipage}[t]{\dimexpr\linewidth-\wd\grafiekbox-3mm\relax}%
}{%
  \end{minipage}\hfill
  \begin{minipage}[t]{\wd\grafiekbox}\vspace{-2mm}\usebox{\grafiekbox}\end{minipage}\par}


% Getallenverzamelingen
\providecommand{\N}{\mathbb{N}}
\providecommand{\Z}{\mathbb{Z}}
\providecommand{\Q}{\mathbb{Q}}
\providecommand{\R}{\mathbb{R}}

% Logische symbolen (zoals in Logica en bewijzen van de 3de graad):
% \niet (negatie), \en (conjunctie), \of (disjunctie), \dan (implicatie), \asa (equivalentie)
\newcommand{\niet}{\neg}
\newcommand{\en}{\wedge}
\newcommand{\of}{\vee}
\newcommand{\dan}{\Rightarrow}
\newcommand{\asa}{\Leftrightarrow}

% WAARHEIDSTABEL: raster in de kleur van het hoofdstuk, alles gecentreerd.
% \begin{waarheidstabel}{5} kop & ... \\ \hline W & O & ... \\ \hline \end{waarheidstabel}
% Het argument is het aantal kolommen. De eerste rij (de kop) is licht getint.
% \res{W} zet een cel van de kolom die je zoekt in een tint, zodat die kolom opvalt.
\newenvironment{waarheidstabel}[1]{%
  \begingroup\arrayrulecolor{accent}\renewcommand{\arraystretch}{1.18}%
  \setlength{\tabcolsep}{2.4mm}\begin{tabular}{|*{#1}{c|}}\hline\rowcolor{accent!10}}%
  {\end{tabular}\endgroup}
% Bredere kolommen, voor lange koppen: \begin{waarheidstabelbreed}{breedte}{aantal kolommen}
\newenvironment{waarheidstabelbreed}[2]{%
  \begingroup\arrayrulecolor{accent}\renewcommand{\arraystretch}{1.18}%
  \begin{tabular}{|*{#2}{>{\centering\arraybackslash}m{#1}|}}\hline\rowcolor{accent!10}}%
  {\end{tabular}\endgroup}
\newcommand{\res}[1]{\cellcolor{accent!22}\textbf{#1}}

% Deelbaarheid: a \deelt b
\newcommand{\deelt}{\mid}
